\documentclass[11pt,a4paper]{article}
\usepackage[utf8]{inputenc}
\usepackage[T1]{fontenc}
\usepackage{amsmath,amssymb,amsthm}
\usepackage{booktabs,array,longtable,caption,placeins,microtype,xcolor}
\usepackage{pgfplots}\pgfplotsset{compat=1.18}
\usepackage[margin=2.5cm]{geometry}
\usepackage[colorlinks=true,linkcolor=blue!60!black,citecolor=blue!60!black,urlcolor=blue!60!black]{hyperref}
\captionsetup{font=small,labelfont=bf}
\IfFileExists{data/wp0.tex}{% generated by scripts/gen-wp0.py
\providecommand{\wpDfone}{0.7507908137}
\providecommand{\wpC}{1.372813462821}
\providecommand{\wpBeta}{1.008}
\providecommand{\wpBetaThree}{0.779}
\providecommand{\wpBetaNuis}{0.991}
\providecommand{\wpBetaNuisThree}{0.998}
\providecommand{\wpRtwo}{0.55}
\providecommand{\wpRtwoThree}{0.34}
\providecommand{\wpBetaZeta}{1.006}
\providecommand{\wpBetaL}{1.020}
\providecommand{\wpRmsResid}{0.528}
\providecommand{\wpRmsLeft}{0.291}
\providecommand{\wpNZ}{452}
\providecommand{\wpNL}{617}
}{}

\newtheorem{theorem}{Theorem}
\newtheorem{proposition}[theorem]{Proposition}
\newtheorem{lemma}[theorem]{Lemma}
\newtheorem{conjecture}[theorem]{Conjecture}
\newtheorem{definition}[theorem]{Definition}
\newtheorem{hypothesis}[theorem]{Hypothesis}
\theoremstyle{remark}\newtheorem{remark}[theorem]{Remark}
\newcommand{\Z}{\mathbb{Z}}\newcommand{\Q}{\mathbb{Q}}\newcommand{\F}{\mathbb{F}}\newcommand{\C}{\mathbb{C}}
\newcommand{\Sig}{\mathfrak S}\newcommand{\Off}{\mathrm{Off}}\newcommand{\Frob}{\operatorname{Frob}}
\newcommand{\Sym}{\operatorname{Sym}}\newcommand{\Res}{\operatorname*{Res}}
% status tags for the draft; remove before submission
\newcommand{\status}[1]{\textcolor{red!70!black}{\small\sffamily[#1]}}

\title{The pair singular series of a polynomial II:\\ zeros of Dedekind zeta functions in the second
moment of prime values}
\author{Jasper Reichardt \and Claude Fable 5.1}
\date{Skeleton, \today}

\begin{document}
\maketitle

\begin{abstract}
\status{draft; numbers to be filled from WP0--WP3}
Let $f\in\Z[t]$ be irreducible without fixed prime divisor, $K=\Q[t]/(f)$, and $S_f(h)$ its pair
singular series. In part I we proved that the diagonal part of $\sum_{h\le H}(S_f(h)-C(f)^2)$ has
Dirichlet series $D_f(s)=\zeta_K(s+1)E_f(s)$. Here we determine the analytic structure of $E_f$
completely: it is an infinite product of Artin $L$-functions of virtual representations of the Galois
group, the first of which is $L(2s+2,\Sym^2V_f)^{-1}=\bigl(\zeta_K(2s+2)\prod_O\zeta_{K_O}(2s+2)\bigr)^{-1}$
with $K_O$ the fields cut out by the unordered pairs of roots; $D_f$ is meromorphic in $\Re s>-1$, has
real poles at $s=-1+1/k$, and $\Re s=-1$ is a natural boundary. Consequently the Riesz means of the
diagonal satisfy an explicit formula, unconditional for $m\ge2$, whose oscillating part is indexed by the
zeros of these Dedekind zeta functions, with terms $x^{m-1+\rho/2}$, and whose smooth part contains
$x^{m-2/3}$. For $f=t^2+1$ every coefficient is computed exactly and the formula is verified against
the exact Riesz means to $x=4\cdot10^7$: \status{WP0 result}. We prove $\Omega_\pm(x^{m-3/4})$
unconditionally, $O(x^{m-7/8+\varepsilon})$ after the real terms under GRH for $\zeta_K$, and give the
function-field version, where the formula is an exact finite sum. Assuming the Hardy--Littlewood pair
conjecture with a power-saving error, these results describe the second-order term of the variance of
the number of prime values of $f$; we explain precisely why the zero-induced fluctuations lie below the
scale at which that transfer is valid, so that no statement about prime values themselves and the zeros
of $\zeta_K$ is claimed.
\end{abstract}

\section{Introduction}\label{sec:intro}

\paragraph{Recollection of part I.} \status{copy the definitions: $\omega_f$, $C(f)$, $S_f$, $b(q)$, $W_f$,
$a_f(d)=W_f(d)\omega_f(d)$, $A_f=a_f*1$, Theorems 2--4 of part I.} The identity
\[
\sum_{h\le H}\bigl(S_f(h)-C^2\bigr)=-C^2\sum_{d\ge2}a_f(d)\Bigl\{\frac Hd\Bigr\}+C^2\Off_f(H)
\]
splits the sum into a diagonal, generated by $D_f(s)=\sum_da_f(d)d^{-s}=\zeta_K(s+1)E_f(s)$ with
$\Res_{s=0}D_f=1/C(f)$, and an off-diagonal over pairs of distinct roots. Conjecture~1 of part I,
$\sum_{h\le H}(S_f(h)-C^2)=-\tfrac12C\log H+A_f+o(1)$, is equivalent (Ces\`aro form) to
$\Off^\ast_f(H)=o(H\log H)$.

\paragraph{The linear case as template.} For $f(t)=t$, Goldston and
Suriajaya~\cite{GoldstonSuriajaya2021} proved that the Riesz means $\sum_{k\le x}(x-k)^m\Sig(k)$, $m\ge2$,
equal $x^{m+1}/(m+1)-\tfrac12x^m(\log x-H_m+\gamma+\log2\pi)+x^{m-1}\sum_{|\gamma|\le U}a(\rho)x^{\rho/2}+O(x^{m-1+\varepsilon})$,
the zeros of $\zeta$ entering through the factor $1/\zeta(2s+2)$ of the generating series; and
$\Omega_\pm(x^{m-3/4})$. \status{cite also Vaughan 2001, Friedlander--Goldston 1995, GS JNT 2021.}

\paragraph{Results.} \status{one paragraph per theorem, in the order of the sections.}

\paragraph{What is conditional on what.} \status{Reuse part I's paragraph (A)/(B)/(C) and add:}
(D)~Everything in Sections~\ref{sec:structure}--\ref{sec:omega} is unconditional and concerns the
arithmetic function $A_f$; the Riemann hypothesis for $\zeta_K$ enters only in the sharp form of the
error term. (E)~Section~\ref{sec:primes} is conditional on the Hardy--Littlewood pair conjecture for $f$
with an error term $O(T^{1/2+\varepsilon})$, as in Montgomery--Soundararajan~\cite{MontgomerySoundararajan2004};
the zero terms of the explicit formula are of relative size $T^{-1/4}$, \emph{below} that tolerance.
Unconditionally, the low-lying zeros of $\zeta_K$ govern the second-order fluctuations of the Riesz
means of $S_f$; under a Bateman--Horn pair conjecture with power saving these means control the
variance of prime values of $f$ only at a scale above the zero-induced fluctuations; and under GRH the
variance is in turn equivalent, by the Goldston--Montgomery theorem~\cite{GoldstonMontgomery1987} in
the form of Bui, Keating and Smith~\cite{BuiKeatingSmith2016}, to pair correlation of the zeros of
$\zeta_K$. We claim no unconditional link between prime values of $f$ and zeros of $\zeta_K$.

\paragraph{Relation to the literature.} \status{from LITERATURE.md \S0: Kuperberg--Rodgers--Roditty-Gershon
(prime elements of $K$, $\zeta_K(s+1)$ in a singular-series average, no explicit formula);
Bhowmik--Halupczok--Matsumoto--Suzuki and Goldston--Suriajaya 2023 (multi-$L$ explicit formulas in the
Goldbach setting); Hall--Keating--Roditty-Gershon and Hochfilzer (variances governed by higher-degree
$L$-functions over $\F_q[t]$); no prior appearance of $\Sym^2V_f$ or of a polynomial singular-series
explicit formula.}

\section{The analytic structure of $D_f$}\label{sec:structure}

Throughout, $f\in\Z[t]$ is irreducible of degree $n$, without fixed prime divisor, $G=\operatorname{Gal}(f)$,
$V=V_f$ the permutation representation of $G$ on the roots, $K=\Q[t]/(f)$, and for a prime $p$ unramified
in the splitting field, $g_p=\Frob_p$ (a conjugacy class) with $\omega_f(p)=\chi_V(g_p)$ fixed points. We
write $w=s+1$, $u=u_p=p^{-1}$, $v=v_p=p^{-w}$. By part I, Theorem~3, $D_f(s)=\prod_pF_p(s)$ with
\begin{equation}\label{eq:Fp}
F_p(s)=P_p+p\,b(p)\,\omega_f(p)\,p^{-s}=\frac{1-2\omega_f(p)u+\omega_f(p)v}{(1-\omega_f(p)u)^2}\qquad(\omega_f(p)<p),
\end{equation}
as one checks by writing $P_p=1-\omega^2u^2/(1-\omega u)^2$ and $p\,b(p)\omega p^{-s}=\omega v/(1-\omega u)^2$.
The point of this section is that $F_p(s)\to1+\omega_f(p)v$ as $u\to0$ with $v$ fixed, and that
$1+\omega_f(p)v=1+\chi_V(g_p)v$ is \emph{not} a finite product of Euler factors of Artin $L$-functions
unless $\chi_V\le1$, i.e.\ unless $f$ is linear.

\subsection{The pair fields}

\begin{lemma}[pair fields]\label{lem:pairs}
$\Sym^2V\cong V\oplus V^{(2)}$, where $V^{(2)}$ is the permutation representation of $G$ on unordered pairs
of distinct roots. Consequently
\[
L(s,\Sym^2V)=\zeta_K(s)\prod_{O}\zeta_{K_O}(s),
\]
the product over the $G$-orbits $O$ on unordered pairs, $K_O$ being the fixed field of the stabiliser of a
pair in $O$. If $G$ is $2$-transitive there is one orbit, and $K^{(2)}=\Q(\alpha+\beta,\alpha\beta)$ has
degree $n(n-1)/2$.
\end{lemma}
\begin{proof}
$\Sym^2V$ has the basis $e_\alpha e_\beta$ ($\alpha\le\beta$ in any fixed order of the roots), which $G$
permutes; the vectors $e_\alpha^2$ span a copy of $V$ and the vectors $e_\alpha e_\beta$, $\alpha\ne\beta$,
span $V^{(2)}$. A permutation representation $\operatorname{Ind}_H^G\mathbf 1$ has Artin $L$-function
$\zeta_{L^H}$, $L$ the splitting field; the stabiliser of the pair $\{\alpha,\beta\}$ fixes
$\alpha+\beta$ and $\alpha\beta$, and $\Q(\alpha+\beta,\alpha\beta)$ has the right degree when $G$ is
$2$-transitive.
\end{proof}

For quadratics, $\Sym^2V=\mathbf 1\oplus\mathbf 1\oplus\chi_D$ and $L(s,\Sym^2V)=\zeta(s)^2L(s,\chi_D)$; for an
$S_3$-cubic $K^{(2)}\cong K$ and $L(s,\Sym^2V)=\zeta_K(s)^2$; for an $S_4$- or $A_4$-quartic $K^{(2)}$ is the
resolvent sextic field.

\subsection{The plethystic exponents}

\begin{definition}\label{def:psi}
For $N\ge1$ and $g\in G$ put
\[
\Psi_N(g)=\frac1N\sum_{j\mid N}\mu(j)\,(-1)^{N/j+1}\,\chi_V(g^j)^{N/j}.
\]
\end{definition}

\begin{lemma}\label{lem:psi}
(i) Each $\Psi_N$ is a virtual character of $G$, and $\Psi_1=\chi_V$.
(ii) For every $g\in G$, as formal power series in $v$,
\[
1+\chi_V(g)\,v=\prod_{N\ge1}\exp\Bigl(\sum_{j\ge1}\frac{\Psi_N(g^j)}{j}v^{Nj}\Bigr)
=\prod_{N\ge1}\det\bigl(1-g\,v^N\mid\Psi_N\bigr)^{-1},
\]
where for a virtual character $\Psi=A-B$ we write $\det(1-gX\mid\Psi)^{-1}=\det(1-gX\mid A)^{-1}\det(1-gX\mid B)$.
(iii) $|\Psi_N(g)|\le 2n^N/N$.
\end{lemma}
\begin{proof}
(ii) Taking logarithms, the claim is $\sum_{M\ge1}c_M(g)v^M=\sum_{N,j}\Psi_N(g^j)v^{Nj}/j$ with
$c_M(g)=(-1)^{M+1}\chi_V(g)^M/M$, i.e.\ $c_M(g)=\sum_{j\mid M}\Psi_{M/j}(g^j)/j$ for all $M$ and $g$. Insert the
definition: $\sum_{j\mid M}\frac1j\cdot\frac{j}{M}\sum_{i\mid M/j}\mu(i)(-1)^{M/(ij)+1}\chi_V(g^{ij})^{M/(ij)}
=\frac1M\sum_{k\mid M}(-1)^{M/k+1}\chi_V(g^k)^{M/k}\sum_{i\mid k}\mu(i)=c_M(g)$, since only $k=1$ survives.
(i) Write $R(g,v)=1+\chi_V(g)v$; its coefficients lie in the representation ring $R(G)$. Suppose
$\Psi_1,\dots,\Psi_{K-1}$ are virtual characters. The series
$R(g,v)\prod_{N<K}\det(1-gv^N\mid\Psi_N)$ has coefficients in $R(G)$ (products of the series
$\det(1-gX\mid A)^{\pm1}$, whose coefficients are $\pm$ the characters of exterior or symmetric powers of
$A$), and by (ii) it equals $1+\Psi_K(g)v^K+O(v^{K+1})$; so $\Psi_K\in R(G)$. The case $N=1$ is the
definition. (iii) $|\chi_V(g^j)|\le n$, and $\sum_{j\mid N}n^{N/j}\le n^N+N n^{N/2}\le2n^N$ for $n\ge2$; for $n=1$ the claim is trivial.
\end{proof}

For $f$ linear, $\chi_V=1$ and $1+v=(1-v^2)/(1-v)$: $\Psi_1=1$, $\Psi_2=-1$, $\Psi_N=0$ for $N\ge3$, which is
the factor $\zeta(s+1)/\zeta(2s+2)$ of the twin-prime case~\cite{GoldstonSuriajaya2021}. For $n\ge2$
infinitely many $\Psi_N$ are non-zero, because $1+\chi_V(g)v$ with $\chi_V(g)\ge2$ is not a finite product
of cyclotomic polynomials in $v$.

\subsection{The structure theorem}

For a virtual character $\Psi$ of $G$ let $L(s,\Psi)$ be its Artin $L$-function (a quotient of two
Artin $L$-functions of genuine representations, hence meromorphic in $\C$), with local factor
$L_p(s,\Psi)=\det(1-g_pp^{-s}\mid\Psi)^{-1}$ at unramified $p$; and let $S$ be the finite set of primes
ramified in the splitting field or dividing the leading coefficient of $f$.

\begin{theorem}[structure of $D_f$]\label{thm:structure}
For every $K\ge1$,
\begin{equation}\label{eq:structure}
D_f(s)=\prod_{N=1}^{K}L\bigl(N(s+1),\Psi_N\bigr)\cdot M_{f,K}(s),\qquad
M_{f,K}(s)=\prod_pF_p(s)\prod_{N=1}^{K}L_p\bigl(N(s+1),\Psi_N\bigr)^{-1},
\end{equation}
where the Euler product $M_{f,K}$ converges absolutely and locally uniformly in $\Re s>-1+1/(K+1)$, so
that $M_{f,K}$ is holomorphic there. In particular $L(s+1,\Psi_1)=\zeta_K(s+1)$, so that
$E_f=\prod_{N=2}^KL(Nw,\Psi_N)\,M_{f,K}$; $D_f$ is meromorphic in $\Re s>-1$; and in
$\Re s>-1+1/(K+1)$ its poles are among the poles of $\prod_{N\le K}L(Nw,\Psi_N)$, while its zeros that are not
zeros of that product are exactly the zeros of the local factors $F_p$, $p\notin S$, together with those
of the finitely many local factors at $p\in S$.
\end{theorem}
\begin{proof}
Fix $K$ and a compact set in $\Re w>1/(K+1)$, on which $\Re w\ge\eta>1/(K+1)$. For $p\notin S$ the local
factor of $M_{f,K}$ is, by Lemma~\ref{lem:psi}(ii),
\[
M_p(s)=\frac{F_p(s)}{1+\omega v}\cdot\frac{1+\omega v}{\prod_{N\le K}\det(1-g_pv^N\mid\Psi_N)^{-1}}
=\frac{1-2\omega u/(1+\omega v)}{(1-\omega u)^2}\cdot\exp\Bigl(-\sum_{M>K}d_M(g_p)v^M\Bigr),
\]
with $\omega=\omega_f(p)$ and $d_M(g)=\sum_{N\mid M,\ N>K}\Psi_N(g^{M/N})N/M$, so $|d_M|\le2n^M\tau(M)$ by
Lemma~\ref{lem:psi}(iii). For $p$ so large that $n|v|\le n p^{-\eta}\le\tfrac12$ and $\omega u\le\tfrac14$
we get $\bigl|\sum_{M>K}d_Mv^M\bigr|\le\sum_{M>K}2\tau(M)(n|v|)^M\ll_K(n|v|)^{K+1}\ll p^{-(K+1)\eta}$, and
$\frac{1-2\omega u/(1+\omega v)}{(1-\omega u)^2}=1+\frac{2\omega^2uv}{1+\omega v}+O(u^2)=1+O(p^{-1-\eta})+O(p^{-2})$.
Since $(K+1)\eta>1$, $\sum_p|M_p(s)-1|$ converges uniformly on the compact set, which gives absolute and
locally uniform convergence of $\prod_{p\notin S,\,p>p_0}M_p$; the finitely many remaining factors are
rational functions of $p^{-s}$ whose only possible poles are at $|p^{-Nw}|=1$, i.e.\ on $\Re w=0$. A
uniformly convergent product of non-vanishing holomorphic factors is holomorphic and non-vanishing; hence
$M_{f,K}$ vanishes exactly where a local factor does, and $M_p(s)=0$ iff $F_p(s)=0$ (the other factors of
$M_p$ are $\det(1-g_pv^N\mid\Psi_N)$, non-zero for $\Re w>0$). Letting $K\to\infty$ gives the
continuation to $\Re s>-1$, the finitely many $L(Nw,\Psi_N)$ being meromorphic by Brauer's theorem.
\end{proof}

\begin{remark}
$M_{f,K}$ is \emph{not} $1+O(p^{-2})$ at the level of individual primes: the mixed term $2\omega^2uv/(1+\omega v)$
is of size $p^{-1-\Re w}$. In the numerical work these terms are extracted as well
($[\zeta(w+1)L(w+1,\chi_D)]^4$, $[\zeta(2w+1)L(2w+1,\chi_D)]^{-8}$, \dots\ for quadratics) so that the
remaining product converges like $\sum p^{-2}$; see \texttt{explicit-diag.py}.
\end{remark}

\subsection{Quadratics: the exponents in closed form}

For a quadratic $f$ of discriminant $D$, $G=\{e,\tau\}$, $\chi_V=\mathbf 1+\chi_D$, and every virtual
character is $a\mathbf 1+b\chi_D$ with $a,b\in\Z$. Write $\Psi_N=a_N\mathbf 1+b_N\chi_D$, so that
$L(Nw,\Psi_N)=\zeta(Nw)^{a_N}L(Nw,\chi_D)^{b_N}$, and let $M_N(x)=\frac1N\sum_{j\mid N}\mu(N/j)x^j$ be the
necklace polynomial.

\begin{proposition}\label{prop:exponents}
For every $N\ge1$,
\[
a_N+b_N=\Psi_N(e)=-M_N(-2),\qquad
a_N-b_N=\Psi_N(\tau)=\frac1N\sum_{\substack{j\mid N\\ 2\mid j}}\mu(j)(-1)^{N/j+1}2^{N/j}.
\]
In particular $a_N=b_N=\tfrac12M_N(2)$ for odd $N$, and for even $N$
$a_N+b_N=-M_N(-2)<0$; the first values are those of Table~\ref{tab:exponents}, and $a_1=b_1=1$, i.e.\
$L(w,\Psi_1)=\zeta_K(w)$; $\Psi_2=-\Sym^2V$; $a_3=b_3=1$, so that $\zeta(3w)L(3w,\chi_D)$ appears in the
numerator and $D_f$ has a simple pole at $s=-\tfrac23$; and $a_N>0$ for all odd $N$, giving a pole of order
$a_N=\tfrac12M_N(2)$ at $s=-1+1/N$.
\end{proposition}
\begin{proof}
$\chi_V(e^j)=2$ and $\chi_V(\tau^j)=2$ or $0$ according as $j$ is even or odd; insert into
Definition~\ref{def:psi}, using $(-1)^{k+1}2^k=-(-2)^k$ for the first formula. For odd $N$ there is no even
divisor, so $a_N=b_N=-\tfrac12M_N(-2)=\tfrac12M_N(2)$; $M_N(2)>0$ counts binary Lyndon words.
\end{proof}

\subsection{The natural boundary}

\begin{theorem}[natural boundary]\label{thm:boundary}
Let $f$ be irreducible of degree $n\ge1$ without fixed prime divisor. Then $D_f$ has infinitely many zeros
in $-1<\Re s<0$, accumulating at every point of the line $\Re s=-1$, and $\Re s=-1$ is a natural boundary
of $D_f$ and of $\sum_nA_f(n)n^{-s}=\zeta(s)D_f(s)$. The same holds for $f(t)=t$, i.e.\ for the
generating Dirichlet series of the twin-prime singular series.
\end{theorem}
\begin{proof}
By~\eqref{eq:Fp}, for $p\notin S$ with $\omega=\omega_f(p)\ge1$, $F_p(s)=0$ iff $\omega v=-(1-2\omega u)$, i.e.\
iff $p^{-w}=-(1-2\omega/p)/\omega$. These are the points
\[
s_{p,j}=-1+\frac{\log\omega-\log(1-2\omega/p)}{\log p}+\frac{(2j+1)\pi i}{\log p},\qquad j\in\Z,
\]
with $\Re s_{p,j}\to-1^+$ as $p\to\infty$ (for $\omega=1$, $\Re s_{p,j}=-1+2/(p\log p)+O(p^{-2})$). By the
Chebotarev density theorem there are infinitely many $p$ with $\omega_f(p)\ge1$ (indeed with
$\omega_f(p)=n$). Given $t_0\in\mathbb R$ and $\varepsilon>0$, choose such a $p$ with $\log p>\pi/\varepsilon$
and $|\Re s_{p,j}+1|<\varepsilon$; the spacing of the $\Im s_{p,j}$ is $2\pi/\log p<2\varepsilon$, so some
$s_{p,j}$ lies within $2\varepsilon$ of $-1+it_0$. By Theorem~\ref{thm:structure} with $K$ large, in a
neighbourhood of $s_{p,j}$ we have $D_f=\Lambda_K\cdot M_{f,K}$ with $\Lambda_K=\prod_{N\le K}L(Nw,\Psi_N)$
meromorphic and $M_{f,K}(s_{p,j})=0$; so $s_{p,j}$ is a zero of $D_f$ unless it is a pole of $\Lambda_K$.
The poles of $\Lambda_K$ form a discrete set, so all but finitely many of the zeros $s_{p,j}$ in any bounded
region are zeros of $D_f$. Now suppose $D_f$ continued meromorphically to a disc $B$ around $-1+it_0$.
Zeros of a meromorphic function accumulate only at poles, so $-1+it_0$ would be a pole; but in a
punctured neighbourhood of a pole $|D_f|$ is large, whereas zeros $s_{p,j}$ lie arbitrarily close to
$-1+it_0$: a contradiction. For $\zeta(s)D_f(s)$ the same argument applies, since $\zeta(s)\ne0$ on $\Re s=-1$. For $f(t)=t$, $F_p=(1-2u+v)/(1-u)^2$ vanishes at $p^{-w}=-(1-2/p)$ and the argument is
identical.
\end{proof}

\begin{remark}
For $f(t)=t$ the pure-$v$ part of the local factor is $1+v=(1-v^2)/(1-v)$, so the $L$-function part of
$D_f$ is just $\zeta(w)/\zeta(2w)$ and the contour in~\cite{GoldstonSuriajaya2021} can be moved to
$\Re s=-1+\varepsilon$, crossing the whole critical strip of $\zeta(2s+2)$: this is why the linear case has
an \emph{unconditional} explicit formula with error $O(x^{m-1+\varepsilon})$, and why that error term cannot
be improved by contour shifting. For $n\ge2$ the factors $L(Nw,\Psi_N)$ with $N\ge3$ pile up between
$\Re s=-\tfrac34$ and $\Re s=-1$, and the contour can only be moved a little past the first zero family.
\end{remark}

\section{The explicit formula for the diagonal}\label{sec:explicit}

Let $R_m(x)=\sum_{n\le x}(x-n)^mA_f(n)$, so that for $\Re s>1$
\begin{equation}\label{eq:mellin}
R_m(x)=\frac1{2\pi i}\int_{(2)}G_m(s)x^{s+m}\,ds,\qquad G_m(s)=\frac{m!\,\zeta(s)D_f(s)}{s(s+1)\cdots(s+m)}.
\end{equation}
From now on $f$ is quadratic with discriminant $D$, so that by Theorem~\ref{thm:structure} and
Proposition~\ref{prop:exponents}, for every $K\ge4$,
\begin{equation}\label{eq:Dquad}
D_f(s)=\frac{\zeta(w)L(w,\chi_D)\,\zeta(3w)L(3w,\chi_D)}{\zeta(2w)^2L(2w,\chi_D)}\;
\prod_{N=4}^{K}\zeta(Nw)^{a_N}L(Nw,\chi_D)^{b_N}\;M_{f,K}(s),\qquad w=s+1.
\end{equation}

\subsection{Statement}

\begin{theorem}\label{thm:explicit-quadratic}
Let $f$ be an irreducible quadratic of discriminant $D$ and $m\ge2$. Assume the Riemann hypothesis for
$\zeta$ and for $L(s,\chi_D)$. Then there is $\varepsilon_0\in(0,\tfrac1{100})$ such that for $x\ge2$ and some
$U\in[x^{16},2x^{16}]$,
\begin{equation}\label{eq:explicit}
R_m(x)=\frac{D_f(1)}{m+1}x^{m+1}+x^m\Bigl(-\frac{\log x}{2C(f)}+c_{f,m}\Bigr)+c'_{f,m}\,x^{m-2/3}
+x^{m-1}\sum_{\substack{\rho=\beta+i\gamma\\|\gamma|\le2U}}Q_{m,\rho}(\log x)\,x^{\rho/2}+O(x^{m-3/4-\varepsilon_0}),
\end{equation}
where $\rho$ runs over the zeros of $\zeta(s)^2L(s,\chi_D)$ (all with $\beta=\tfrac12$), $Q_{m,\rho}$ is the
polynomial (of degree less than the order of the pole of $G_m$ at $\rho/2-1$, i.e.\ of degree $0$ at simple
zeros of $L$ and $1$ at simple zeros of $\zeta$) such that $Q_{m,\rho}(\log x)x^{\rho/2-1+m}$ is the residue
of $G_m(s)x^{s+m}$ at $s=\rho/2-1$, and
\[
c'_{f,m}=\Res_{s=-2/3}G_m(s)=\frac{m!\,\zeta(-\tfrac23)\zeta(\tfrac13)L(\tfrac13,\chi_D)L(1,\chi_D)}{3\,\zeta(\tfrac23)^2L(\tfrac23,\chi_D)}
\prod_{N\ge4}\zeta(\tfrac N3)^{a_N}L(\tfrac N3,\chi_D)^{b_N}\frac{M_{f,K}(-\tfrac23)}{(-\tfrac23)(\tfrac13)\cdots(m-\tfrac23)},
\]
$c_{f,m}$ being the constant term of the residue at $s=0$. The implied constant depends on $f$, $m$ and
$\varepsilon_0$.
\end{theorem}

The coefficient of $x^m\log x$ is $-1/(2C(f))$ because $\Res_{s=0}D_f=1/C(f)$ (part I, Theorem~3) and
$\zeta(0)=-\tfrac12$; explicitly $c_{f,m}=-\frac1{2C(f)}\bigl(\log2\pi-H_m\bigr)-\tfrac12\kappa_f$ with
$\kappa_f$ the constant term of the Laurent expansion of $D_f$ at $0$ and $H_m=\sum_{k\le m}1/k$.

\subsection{Growth lemmas}

We use the following consequences of the Riemann hypothesis for $\zeta$ and $L(s,\chi_D)$
(\cite[Chapter~14]{Titchmarsh1986} for $\zeta$; the proofs transfer verbatim to $L(s,\chi_D)$, see
\cite[\S5.7]{IwaniecKowalski2004}). Let $\Lambda\in\{\zeta,L(\cdot,\chi_D)\}$ and $t\ge2$.
\begin{itemize}
\item[(G1)] For every fixed $\sigma>\tfrac12$: $\Lambda(\sigma+it)^{\pm1}\ll t^{\varepsilon}$; for every fixed $\sigma$:
$\Lambda(\sigma+it)\ll t^{\max(0,1/2-\sigma)+\varepsilon}$, and for fixed $\sigma<\tfrac12$:
$\Lambda(\sigma+it)^{-1}\ll t^{\varepsilon}$ (by the functional equation, $|\Lambda(\sigma+it)|\asymp t^{1/2-\sigma}|\Lambda(1-\sigma-it)|$).
\item[(G2)] Unconditionally, for $-1\le\sigma\le2$: $\dfrac{\Lambda'}{\Lambda}(\sigma+it)=\sum_{|\gamma-t|\le1}\dfrac1{\sigma+it-\rho}+O(\log t)$,
the sum over zeros $\rho=\beta+i\gamma$ of $\Lambda$, and the number of zeros with $|\gamma-t|\le1$ is $\le c_0\log t$.
\end{itemize}

\begin{lemma}[horizontal segments]\label{lem:horizontal}
Assume RH for $\zeta$ and $L(s,\chi_D)$. For every $\delta>0$ there is $\varepsilon_1(\delta)>0$ such that
for all $T\ge T_0(\delta)$ and $\varepsilon_0\le\varepsilon_1$, there exists $U\in[T,2T]$ with
\[
\bigl|\zeta(\sigma+2iU)^{-1}\bigr|+\bigl|L(\sigma+2iU,\chi_D)^{-1}\bigr|\le T^{\delta}\qquad\text{uniformly for }\tfrac12-2\varepsilon_0\le\sigma\le2 .
\]
\end{lemma}
\begin{proof}
Let $\Lambda$ be one of the two functions, $\sigma_1=\tfrac12+\varepsilon_0$, and $Z=[\tfrac12-2\varepsilon_0,\sigma_1]$
(the ``zone''). For $\sigma\ge\sigma_1$ the claim is (G1) with $\delta/2$ in place of $\varepsilon$. For $\sigma\in Z$ and
$t=2U$ write, by (G2),
\[
\log\frac1{|\Lambda(\sigma+it)|}=\log\frac1{|\Lambda(\sigma_1+it)|}+\Re\int_\sigma^{\sigma_1}\frac{\Lambda'}{\Lambda}(\sigma'+it)\,d\sigma'
\le\frac{\delta}4\log T+\sum_{|\gamma-t|\le1}\int_{Z}\frac{d\sigma'}{|\sigma'+it-\rho|}+3\varepsilon_0C\log t .
\]
Under RH $\beta=\tfrac12$, so $|\sigma'+it-\rho|=\sqrt{(\sigma'-\tfrac12)^2+(t-\gamma)^2}$ and, with
$\theta=|t-\gamma|$, $\int_Z\le\int_{-2\varepsilon_0}^{2\varepsilon_0}\frac{dy}{\sqrt{y^2+\theta^2}}=2\operatorname{arsinh}(2\varepsilon_0/\theta)$,
which is $\le4\varepsilon_0/\theta$ for $\theta\ge2\varepsilon_0$ and $\le2\log(4\varepsilon_0/\theta)+2$ for $\theta<2\varepsilon_0$.
By (G2) the zeros with $2\varepsilon_0\le\theta\le1$ contribute at most $4\varepsilon_0\sum\theta^{-1}\le4\varepsilon_0\log(1/2\varepsilon_0)\,c_0\log t$
(split the range of $\theta$ dyadically: each dyadic block $[\theta_0,2\theta_0]$ contains at most $c_0\log t$ zeros, each contributing at most $\theta_0^{-1}$). The zeros with $\theta<2\varepsilon_0$ contribute $2\Sigma_\Lambda(t)+2\#\{\gamma:|\gamma-t|<2\varepsilon_0\}$
with $\Sigma_\Lambda(t)=\sum_{|\gamma-t|<2\varepsilon_0}\log(4\varepsilon_0/|t-\gamma|)$. Now average over $t\in[2T,4T]$:
each zero contributes at most $\int_{-2\varepsilon_0}^{2\varepsilon_0}\log(4\varepsilon_0/|y|)\,dy=4\varepsilon_0(1+\log2)$, and there are
$\le c_0'T\log T$ zeros with $\gamma\in[2T-1,4T+1]$, so $\frac1{2T}\int_{2T}^{4T}\Sigma_\Lambda(t)\,dt\le 6c_0'\varepsilon_0\log T$
and similarly the average of $\#\{\gamma:|\gamma-t|<2\varepsilon_0\}$ is $\le4c_0'\varepsilon_0\log T$. By Markov's inequality the set
of $t\in[2T,4T]$ with $\Sigma_\zeta(t)+\Sigma_L(t)+\#_\zeta+\#_L\le 100c_0'\varepsilon_0\log T$ has measure $\ge\tfrac{3}{5}\cdot2T$;
pick $U$ with $2U$ in it. Altogether $\log(1/|\Lambda(\sigma+2iU)|)\le(\tfrac\delta4+c_1\varepsilon_0\log(1/\varepsilon_0))\log T$
for $\sigma\in Z$, with $c_1$ absolute, which is $\le\delta\log T$ once $\varepsilon_0\le\varepsilon_1(\delta)$.
\end{proof}

\subsection{Proof of Theorem~\ref{thm:explicit-quadratic}}

Fix $K=15$ (any $K\ge4$ with $(K+1)(\tfrac14-\varepsilon_0)>1$ works), $\delta=\varepsilon_0$ and
$\varepsilon_0\le\min(\varepsilon_1(\varepsilon_0),\tfrac1{100})$, which is possible since $\varepsilon_1(\delta)\gg\delta/\log(1/\delta)$.
Let $U$ be as in Lemma~\ref{lem:horizontal} with $T=x^{16}$, and let $\mathcal R$
be the rectangle with vertices $2\pm iU$ and $-\tfrac34-\varepsilon_0\pm iU$. By~\eqref{eq:mellin} (truncating the
integral at height $U$ costs $\ll x^{m+2}U^{-m}$, since $G_m(2+it)\ll t^{-m-1}$) and the residue theorem,
\[
R_m(x)=\sum_{\text{poles in }\mathcal R}\Res\bigl(G_m(s)x^{s+m}\bigr)+\frac1{2\pi i}\Bigl(\int_{-3/4-\varepsilon_0-iU}^{-3/4-\varepsilon_0+iU}+\int_{\text{horiz.}}\Bigr)G_m(s)x^{s+m}\,ds+O(x^{m+2}U^{-m}).
\]
\emph{Poles.} By~\eqref{eq:Dquad} and Theorem~\ref{thm:structure}, in $\Re s\ge-\tfrac34-\varepsilon_0>-\tfrac45$
the poles of $G_m$ are: $s=1$ (simple, from $\zeta(s)$; $D_f(1)=\sum_da_f(d)/d$), $s=0$ (double: the pole of
$\zeta(w)$ and the factor $1/s$), $s=-\tfrac23$ (simple, from $\zeta(3w)$), and $s=\rho/2-1$ for the zeros $\rho$
of $\zeta(2w)^2L(2w,\chi_D)$ with $|\gamma|\le2U$; the factors $\zeta(Nw)^{a_N}L(Nw,\chi_D)^{b_N}$, $N\ge4$, have
$\Re(Nw)\ge1-4\varepsilon_0>\tfrac12$ there, so under RH they have neither zeros nor poles except
$\zeta(Nw)$ at $Nw=1$, i.e.\ $s=-1+1/N\le-\tfrac34$, outside $\mathcal R$ for $N\ge4$ (and $N=4$ gives a zero
of $D_f$); $M_{f,K}$ is holomorphic in $\Re s>-1+\tfrac1{16}$; and $s=-1,\dots,-m$ are outside $\mathcal R$.
The residues at $1$, $0$, $-\tfrac23$ are the three explicit terms of~\eqref{eq:explicit}, and the residue at
$\rho/2-1$ is $Q_{m,\rho}(\log x)x^{\rho/2-1+m}$ by definition.

\emph{Vertical line.} On $\Re s=-\tfrac34-\varepsilon_0$ (so $\Re w=\tfrac14-\varepsilon_0$, $\Re2w=\tfrac12-2\varepsilon_0$,
$\Re3w=\tfrac34-3\varepsilon_0$, $\Re Nw\ge1-4\varepsilon_0$ for $N\ge4$), (G1) gives, for $|t|\ge2$,
$\zeta(s)\ll|t|^{5/4+\varepsilon_0+\varepsilon}$, $\zeta(w)L(w,\chi_D)\ll|t|^{1/2+2\varepsilon_0+\varepsilon}$,
$\zeta(2w)^{-2}L(2w,\chi_D)^{-1}\ll|t|^{\varepsilon}$, $\zeta(3w)L(3w,\chi_D)\ll|t|^\varepsilon$,
$\prod_{N\ge4}\ll|t|^{\varepsilon}$, and $M_{f,K}\ll1$; hence $G_m(s)\ll|t|^{7/4+3\varepsilon_0+\varepsilon-(m+1)}$, and
for $m\ge2$ the exponent is $<-1$. The vertical integral is therefore $\ll x^{m-3/4-\varepsilon_0}$.

\emph{Horizontal segments.} On $\Im s=\pm U$, $-\tfrac34-\varepsilon_0\le\sigma\le2$: (G1) bounds $\zeta(s)$,
$\zeta(w)L(w)$, $\zeta(3w)L(3w)$ and the $N\ge4$ factors by $U^{7/4+3\varepsilon_0+\varepsilon}$ (worst case at the left end),
Lemma~\ref{lem:horizontal} bounds $\zeta(2w)^{-2}L(2w,\chi_D)^{-1}$ by $U^{3\varepsilon_0}$, and $M_{f,K}\ll1$; the kernel
is $\ll U^{-m-1}$ and $|x^{s+m}|\le x^{m+2}$. The two segments contribute
$\ll x^{m+2}U^{7/4+6\varepsilon_0+\varepsilon-m-1}\le x^{m+2}U^{-1/4+7\varepsilon_0}$ for $m\ge2$, and the truncation of
the original integral at height $U$ contributes $\ll x^{m+2}U^{-m}$; with $U\ge x^{16}$ and $\varepsilon_0\le\tfrac1{100}$
both are $\ll x^{m-1}$.

Collecting, $R_m(x)$ equals the residue sum plus $O(x^{m-3/4-\varepsilon_0})$, which is~\eqref{eq:explicit}. \qed

\begin{remark}[on the zero sum]
The sum over $|\gamma|\le2U$ is not absolutely convergent as $U\to\infty$ with the polynomial growth of the
coefficients; as in \cite[Theorem~3]{GoldstonSuriajaya2021}, simplicity of the zeros and a bound
$\sum_{|\gamma|\le T}|\zeta'(\rho)|^{-2}\ll T$ (Gonek--Hejhal) would make $x^{m-3/4}\sum_\gamma Q_{m,\rho}x^{i\gamma/2}$
absolutely convergent for $m\ge2$. Numerically the coefficients decay like $|\gamma|^{-5/4}$ for $m=2$.
\end{remark}

\subsection{$f=t^2+1$: every coefficient}\label{sec:t2p1}

For $f=t^2+1$, $D=-4$, $\chi_D=\chi_{-4}$, $K=\Q(i)$, and every quantity in~\eqref{eq:explicit} is computable:
$D_f(1)=\wpDfone$, $C(f)=\wpC$ (from $L(1,\chi_{-4})=\pi/4$ and $M_{f,K}(0)$; part I quotes
$1.3728134628\dots$), $c_{f,2}=-0.454974899864$, $c_{f,3}=-0.333569719058$, and the coefficients
$Q_{m,\rho}$ for the zeros with $\gamma\le100$ are computed as Taylor coefficients of
$(s-\rho/2+1)^{\mathrm{ord}}G_m(s)$ by trapezoidal Cauchy integrals on circles of radius $0.015$
(\texttt{explicit-diag.py}; $M_{f,K}$ with $K=15$ and primes up to $4\cdot10^6$, with the tails of the
mixed terms handled by the prime number theorem). \status{Table of $Q_{2,\rho}$ for the first $8$ zeros of
$L(\chi_{-4})$ and $6$ of $\zeta$; generate from the cached residues.}

\subsection{Numerical verification}\label{sec:num}

\IfFileExists{data/explicit-diag-2.dat}{%
\begin{figure}[htbp]\centering
\begin{tikzpicture}
\begin{axis}[width=0.95\textwidth,height=6cm,xlabel={$u=\log x$},ylabel={$/x^{m-3/4}$},legend pos=north west,
  legend style={font=\scriptsize},title={$f=t^2+1$, $m=2$: data minus exact smooth terms (grey) and predicted zero sum (blue)}]
\addplot[gray,thin] table[x=u,y=resid,each nth point=2] {data/explicit-diag-2.dat};
\addplot[blue!70!black,thin] table[x=u,y=pred,each nth point=2] {data/explicit-diag-2.dat};
\legend{residual,prediction}
\end{axis}
\end{tikzpicture}
\caption{The residual of the diagonal Riesz mean of order $2$ of $t^2+1$ after subtracting the exact terms at
$s=1$, $s=0$ and the real poles $s=-2/3,-4/5,-6/7$, against the sum over the zeros of $L(s,\chi_{-4})$ and of
$\zeta$ with $\gamma\le100$. Regression coefficient $\beta=\wpBeta$ ($\wpBetaNuis$ with a smooth nuisance
basis); per family $\beta_\zeta=\wpBetaZeta$, $\beta_L=\wpBetaL$.}
\label{fig:wp0}
\end{figure}}{}

The raw Riesz means $R_m(x)$, $m=1,2,3$, were computed exactly (compensated summation, $A_f(n)$ from
$\omega_f(p)=1+\chi_{-4}(p)$) at $4096$ points equally spaced in $u=\log x\in[\log10^4,\log4\cdot10^7]$. From
$R_2$ we subtract the exact terms at $s=1$, $s=0$, $-\tfrac23$, $-\tfrac45$ (order $3$) and $-\tfrac67$ (order
$9$); the only fitted quantity is a correction to the $x^{m+1}$ coefficient, which comes out at
$-1.0\cdot10^{-12}$ relative to $D_f(1)/(m+1)$ and is itself a check of $D_f(1)$ to twelve digits. The residual,
divided by $x^{5/4}$, has root mean square $\wpRmsResid$. The predicted zero sum (zeros of $L(s,\chi_{-4})$ and of $\zeta$ with $\gamma\le100$) has root mean square $0.27$. Regressing the residual on the prediction
gives $\beta=\wpBeta$ (Figure~\ref{fig:wp0}); with a smooth nuisance basis $x^{m-8/9}(\log x)^j$, $j\le3$,
absorbing the dropped real poles, $\beta=\wpBetaNuis$ with $R^2=\wpRtwo$ and a remaining root mean square of
$\wpRmsLeft$; the two families separately give $\beta_\zeta=\wpBetaZeta$ and $\beta_L=\wpBetaL$. Controls in
which all ordinates are shifted by $\delta\in[-0.6,0.6]$ give $\beta$ between $-0.94$ and $0.68$ and negative
$R^2$. For $m=3$ the zero terms (rms $0.11$) lie below the smooth leftover (rms $0.6$); with the nuisance basis
$\beta=\wpBetaNuisThree$. The remaining rms $0.29$ is accounted for by the truncation at $\gamma\le100$
(the coefficients decay only like $\gamma^{-5/4}$) and the $O(x^{m-3/4-\varepsilon_0})$ term. We regard this as
a confirmation of Theorem~\ref{thm:explicit-quadratic} with the predicted amplitudes and phases to about
$1\%$; it does \emph{not} resolve individual ordinates closer than $4\pi/8.3\approx1.5$. \status{say what
happened when the contour was pushed to $-7/8$: the second family has rms $13$ and $92$ and $\beta=0.001$.}

\section{$\Omega$-results and conditional bounds}\label{sec:omega}

\begin{theorem}\label{thm:omega}
Let $f$ be an irreducible quadratic and $m\ge1$, and let
$E_m(x)=R_m(x)-\frac{D_f(1)}{m+1}x^{m+1}-x^m\bigl(-\frac{\log x}{2C(f)}+c_{f,m}\bigr)-c'_{f,m}x^{m-2/3}$.
Suppose that for some zero $\rho_1=\tfrac12+i\gamma_1$ of $\zeta(s)^2L(s,\chi_D)$ the function $G_m$ has a genuine
pole at $\rho_1/2-1$ (i.e.\ $Q_{m,\rho_1}\ne0$). Then, unconditionally,
\[
E_m(x)=\Omega_\pm\bigl(x^{m-3/4}\bigr).
\]
For $f=t^2+1$ the hypothesis holds with $\gamma_1=6.0209\dots$, the first zero of $L(s,\chi_{-4})$
($|Q_{2,\rho_1}|$ is computed in Section~\ref{sec:t2p1} and is non-zero).
\end{theorem}
\begin{proof}
Suppose $E_m(x)\le\varepsilon x^{m-3/4}$ for all $x\ge x_0$. The function
$F(x)=\varepsilon x^{m-3/4}-E_m(x)$ is then non-negative for $x\ge x_0$, and for $\Re s>1$
\[
\int_{x_0}^\infty F(x)\,x^{-s-m-1}\,dx=\frac{\varepsilon x_0^{-s-1/4}}{s+\tfrac14}-G_m(s)+\Bigl(\text{the Mellin transforms of the subtracted terms}\Bigr)+\bigl(\text{entire}\bigr),
\]
where we used $\int_1^\infty R_m(x)x^{-s-m-1}dx=G_m(s)$ for $\Re s>1$ (Mellin inversion of~\eqref{eq:mellin})
and the fact that $\int_1^{x_0}$ is entire; the Mellin transforms of $x^{m+1}$, $x^m\log x$, $x^m$, $x^{m-2/3}$
are $1/(s-1)$, $1/s^2$, $1/s$, $1/(s+\tfrac23)$ up to constants, and they cancel the poles of $G_m$ at $1$, $0$,
$-\tfrac23$ exactly by the definition of the coefficients. The right-hand side is therefore meromorphic in
$\Re s>-\tfrac45$ (Theorem~\ref{thm:structure}), holomorphic on the real segment $(-\tfrac45,\infty)$ (the
remaining real singularities of $D_f$ in this range are the zeros at $-\tfrac12$ and $-\tfrac34$), and has a
pole at $\rho_1/2-1$, of real part $-\tfrac34$. By Landau's theorem for integrals with non-negative
integrand~\cite[Theorem~1.7]{MontgomeryVaughan2007}, the abscissa of convergence of the left-hand side
is a singularity of the function on the real axis; hence the integral converges for $\Re s>-\tfrac45$, and its
value is holomorphic there, contradicting the pole at $\rho_1/2-1$. Thus $E_m(x)>\varepsilon x^{m-3/4}$ for
arbitrarily large $x$, i.e.\ $E_m=\Omega_+(x^{m-3/4})$; applying the argument to $-E_m$ gives $\Omega_-$.
\end{proof}

\begin{remark}
(i) No positivity of $A_f$ is needed, and no hypothesis beyond the existence of one uncancelled pole on
$\Re s=-\tfrac34$; if $\zeta^2L(\chi_D)$ had a zero with $\beta>\tfrac12$ the same argument gives
$\Omega_\pm(x^{m-1+\beta/2})$, provided that pole is uncancelled. (ii) Under RH for $\zeta$ and
$L(s,\chi_D)$, Theorem~\ref{thm:explicit-quadratic} gives $E_m(x)\ll x^{m-3/4+\varepsilon}$ for $m\ge2$
provided the zero sum is $O(x^{\varepsilon})$ in the normalisation $x^{m-3/4}$, which follows from the
simplicity and Gonek--Hejhal hypotheses of the remark above; so under those hypotheses the exponent
$m-\tfrac34$ is exact. (iii) For $f=t^2+t+41$ the first zero of $L(s,\chi_{-163})$ has $\gamma_1=0.2029$, so the
$\Omega$-oscillation has period $4\pi/\gamma_1\approx62$ in $\log x$: over any computable range the term
looks like a constant multiple of $x^{m-3/4}$. The family $\{L(s,\chi_D)\}$ has symplectic symmetry
\cite{KatzSarnak1999}, so such a low zero is the class-number-one exception, not the rule.
\end{remark}

\section{The off-diagonal}\label{sec:off}

\status{Route A of routes.tex: Prop.~A1 (proved), A2/A3 (to prove) for quadratics. New here: the
off-diagonal Dirichlet series $\sum_dd^{-s}W_f(d)\sum_{s\ne s'}\zeta(s,\langle(s'-s)/d\rangle)$ and the
statement that Conjecture~1 of part I, in analytic form, is its continuation to a pole-free strip left of
$0$. Numerical: $\Off_f(H)$ for $t^2+1$ to $10^7$ and its spectrum. If nothing more: state clearly that the
explicit formula of Section~\ref{sec:explicit} is for the diagonal and that the full sum adds an
$O(H)$-in-Ces\`aro-form term whose fluctuations are those of the fractional parts of the root differences.}

\section{The function field $\F_q[u]$}\label{sec:ff}

\status{Route B. New: with $T=q^{-s}$ the diagonal generating function is a rational function of $T$
(Prop.~B2 of routes.tex + Theorem~\ref{thm:structure}), so the Riesz-mean analogue is an EXACT finite sum
over the inverse roots: the Weil zeros of $L(T,\chi_D)$ and of the extracted factors, and the real
"poles" $T=q^{1/k}$. Write it down for $f=t^2-D$; compare with \texttt{thesis/ff.ts} data. Then
Prop.~B3, $q\to\infty$, for the off-diagonal (Katz/Deligne).}

\section{Consequences for prime values}\label{sec:primes}

\status{Conditional section. (i) Under the HL pair conjecture with error $O(T^{1/2+\varepsilon})$, the
variance formula of part I with the diagonal explicit formula inserted: the polynomial
Montgomery--Soundararajan law plus an explicit oscillatory correction of relative size $T^{-1/4}$, which is
below the error tolerated by the hypothesis --- say so. (ii) Goldston--Montgomery for $\zeta_K$: what
Bui--Keating--Smith give for the degree-$2$ non-primitive $\zeta_K=\zeta L(\chi_D)$ (union of two
independent GUE processes, Murty--Perelli), and why no analogue exists for prime values of a non-linear
$f$ (no Dirichlet series). (iii) What GRH for $\zeta_K$ does say as a theorem: Grenié--Molteni--Perelli
for prime ideals; our $O(x^{m-7/8+\varepsilon})$.}

\section{Discussion and open problems}\label{sec:disc}
\status{(1) Prove Theorem~\ref{thm:explicit-general} in general. (2) The off-diagonal for quadratics
(Route A). (3) Is the natural boundary known for the Goldston--Suriajaya series? (4) Higher moments:
$R_k$ for polynomial tuples and Gaussian fluctuations. (5) The constant $A_f$ in closed form.}

\section*{Data and code}
\texttt{followup/scripts/riesz-raw.ts} (raw Riesz means), \texttt{followup/scripts/explicit-diag.py}
(exponents, constants, residues, comparison), \texttt{followup/data/*}. \status{add WP2 scripts}

\section*{Statement on authorship and the use of AI}
\status{as in part I}

\bibliographystyle{plain}
\bibliography{refs}
\end{document}
